% Generated draft: requires mathematical and source review.
\documentclass[11pt]{article}
\usepackage[margin=25mm]{geometry}
\usepackage{fontspec}
\setmainfont{Latin Modern Roman}
\usepackage{amsmath,amssymb,mathtools,unicode-math}
\setmathfont{Latin Modern Math}
\usepackage{enumitem,needspace,xurl,hyperref}
\setlength{\parindent}{0pt}
\setlength{\parskip}{5pt}
\setlength{\emergencystretch}{3em}
\newcommand{\N}{\mathbb{N}}
\newcommand{\Z}{\mathbb{Z}}
\newcommand{\Q}{\mathbb{Q}}
\newcommand{\R}{\mathbb{R}}
\newcommand{\C}{\mathbb{C}}
\newcommand{\F}{\mathbb{F}}
\newcommand{\A}{\mathbb{A}}
\newcommand{\PP}{\mathbb{P}}
\newcommand{\E}{\mathbb{E}}
\newcommand{\eps}{\varepsilon}
\newcommand{\dd}{\,\mathrm{d}}
\begin{document}
\section*{UnsolvedMath ID: 700310}
\textbf{Source problem code:} catalogue entry
\subsection{Definitions and mathematical statement}
Let \(p\) be a prime and let \(G\) be a finite \(p\)-group. Write \(d=d(G)\) for the least cardinality of a generating set of \(G\). Let \(F_d\) be the free discrete group on \(d\) generators and let \(\widehat F_{d,p}\) be the free pro-\(p\) group on \(d\) topological generators. Define
\[
 r_1(G)=\min\bigl\{|R|:G\cong F_d/\langle\!\langle R\rangle\!\rangle,\ R\subseteq F_d\bigr\},
\]
and
\[
 r_2(G)=\min\bigl\{|R|:G\cong\widehat F_{d,p}/\overline{\langle\!\langle R\rangle\!\rangle},\ R\subseteq\widehat F_{d,p}\bigr\}.
\]
Here \(\langle\!\langle R\rangle\!\rangle\) denotes normal closure, the bar denotes closure in the pro-\(p\) topology, and the second isomorphism is an isomorphism of topological groups, with \(G\) given its finite discrete topology.

For each prime \(p\), define the optimal universal constant
\[
 c(p)=\sup\bigl\{b\in\R:\ r_2(G)\ge b\,d(G)^2\text{ for every finite }p\text{-group }G\bigr\}.
\]
Equivalently,
\[
 c(p)=\inf_{\substack{G\text{ a finite }p\text{-group}\\d(G)>0}}
 \frac{r_2(G)}{d(G)^2}.
\]
The trivial group imposes no restriction on the defining non-strict inequality for \(c(p)\). The source also states the background inequalities \(r_1(G)\ge r_2(G)\) and \(r_2(G)>d(G)^2/4\). The latter cannot literally include the trivial group, for which \(d(G)=r_2(G)=0\); its stated scope therefore requires a nontriviality qualification.

\begin{enumerate}
\item[(a)] Find a finite \(p\)-group \(G\) such that
\[
 r_1(G)>r_2(G).
\]
The source does not specify whether \(p\) may be chosen or whether a witness is requested for every prescribed prime.
\item[(b)] The conjecture is
\[
 \lim_{\substack{p\to\infty\\p\text{ prime}}}c(p)=\frac14.
\]
\end{enumerate}

The source additionally attributes to J. Wisliceny the assertion
\[
 \lim_{d\to\infty}\frac{r_2}{d^2}=\frac14.
\]
This is retained as a supplied literature assertion with unresolved quantifiers: \(r_2\) depends on the group, and the record specifies neither a family of groups, an extremum over groups of generator rank \(d\), nor how the prime varies. Thus this displayed limit does not, as supplied, determine a precise auxiliary theorem.
\subsection{Short English statement}
For finite \(p\)-groups, compare the fewest defining relations on a minimal generating set in discrete and pro-\(p\) presentations. Part (a) requests a strict difference between these numbers. Part (b) conjectures that the optimal universal lower bound for the pro-\(p\) relation number divided by the square of the generator number tends to \(1/4\) as the prime increases. The prime quantifier in part (a) and the scope of the quoted Wisliceny limit remain unresolved.
\subsection{Sources}
\begin{itemize}
\item Supplied local UnsolvedMath registry record: permanent ID \(700310\), exact source problem code \texttt{catalogue entry}, title \emph{Kourovka Notebook Problem 5.26}. The registry labels the problem \emph{open}; this is supplied source metadata, not an independently verified status assessment.
\item H. Koch, proposer, 1976. \emph{The Kourovka Notebook (Unsolved Problems in Group Theory)}, Problem 5.26, originating in the fifth issue (1976); the supplied record identifies the twenty-first issue edition (2026), \texttt{arXiv:1401.0300v48}. The collection metadata names E. I. Khukhro and V. D. Mazurov as editors.
\item J. Wisliceny, \emph{Math. Nachr.} \textbf{102} (1981), 57--78. The source cites this reference for the displayed asymptotic assertion about \(r_2/d^2\), without supplying an article title or the quantifiers needed to interpret that assertion.
\end{itemize}
\end{document}
